\documentclass[11pt]{article}
\usepackage{amsmath,amsthm,amssymb,mathtools}
\usepackage[margin=1in]{geometry}
\usepackage{hyperref}
\usepackage{booktabs}
\usepackage{enumitem}

\newtheorem{theorem}{Theorem}
\newtheorem{lemma}[theorem]{Lemma}
\newtheorem{proposition}[theorem]{Proposition}
\newtheorem{corollary}[theorem]{Corollary}
\newtheorem{conjecture}[theorem]{Conjecture}
\theoremstyle{definition}
\newtheorem{definition}[theorem]{Definition}
\newtheorem{remark}[theorem]{Remark}
\newtheorem{example}[theorem]{Example}

\DeclareMathOperator{\im}{im}
\DeclareMathOperator{\rank}{rank}
\DeclareMathOperator{\supp}{supp}
\newcommand{\Hq}{H_n(q)}
\newcommand{\Vlam}{V^\lambda}
\newcommand{\Vnu}{V^\nu}
\newcommand{\Vmu}{V^\mu}
\newcommand{\Rp}{{R'}}
\newcommand{\Ep}[1]{E^{+}_{#1}}
\newcommand{\Em}[1]{E^{-}_{#1}}
\newcommand{\Om}{\Omega^{(\lambda)}_{\mathrm{tail}}}
\newcommand{\OmU}{\Omega^{(\lambda),[\widehat U]}_{\mathrm{tail}}}
\newcommand{\hatl}{\widehat\lambda}
\newcommand{\Pq}{P_{q}}
\newcommand{\Pm}{P_{-1}}

\title{A deficit law for Strong Pillar~1:\\
removing $u$ tail factors costs exactly $u$ levels of sign-kill\\
\large (the off--hook family $\lambda=(2,2,1^m)$, uniformly in $m$)}
\author{Clio}
\date{29 May 2026 (prove session)}

\begin{document}
\maketitle

\begin{abstract}
The strong per--subset operator vanishing $M(S)\equiv0$ on $\Vlam$ for $|S|<D=\tau(\tau+1)/2$
governs the off--hook order law $\mathrm{ord}_{x=q^2}Z_\lambda=D$. Via the
linear--combination identity of the May~28 reduction paper, the per--subset chain
$M_T(S_T)$ is a signed sum of \emph{factor--removed} tail chains
$\OmU{}$, so per--subset Pillar~1 reduces to the uniform
\emph{Strong Pillar~1} statement on the $\OmU{}$. We prove, uniformly in
$m$, the sharp \emph{deficit form} of Strong Pillar~1 for the off--hook family
$\lambda=(2,2,1^m)$:
\[
   \OmU{}\,\Vlam\ \subseteq\ \bigcap_{j=1}^{\tau-|U|}\Em j\big|_{\Vlam}
   \qquad\text{for every }U\subseteq\mathrm{tail},\ |U|<\tau,
\]
i.e.\ deleting $u$ tail factors lowers the guaranteed depth of sign--kill by exactly $u$.
The engine is an elementary structural fact about the contributor hook $(2,1^m)$: on the
hook seminormal module \emph{every} generator $S_i=T_i+1$ acts as a rank--one operator
whose image is supported on the two basis tableaux with arm value in $\{i,i+1\}$. Hence a
descending product $\prod_{i\in A}S_i$ has image pinned by its leftmost factor
$S_{\max A}$, with support arm value $\ge\max A\ge|A|$; this gives the hook deficit bound
with no induction. The off--hook bound follows by the established $\Rp_n$--branching
(only contributor: the hook $(2,1^m)$; no vanishers): when the leftmost tail generator
$T_{n-1}$ survives, the image lands in $\Ep{n-1}=\Phi(V^{(2,1^m)})$ and pulls through to
the hook bound; when $T_{n-1}$ is deleted, the chain lies in $H_q(S_{n-1})$ and
branches to the larger hook $(2,1^{m+1})$ (closed again by the rank--one fact) and the
smaller off--hook $(2,2,1^{m-1})$ (closed by induction on $m$). As a corollary we obtain
refined per--subset Pillar~1 $M_T(S_T)\Vlam\subseteq\bigcap_{j=1}^{\tau-|S_T|}\Em j$ for
all $|S_T|<\tau$ and all $m$, upgrading the May~28 result (which proved only $2$ of $11$
single--factor positions, at $m=3$ only) to a uniform theorem.
\end{abstract}

\section{Setup}\label{sec:setup}

We use the conventions of the May~15 multi--row meta paper and the May~28 per--subset
reduction paper. $\Hq$ is the Iwahori--Hecke algebra with $T_i^2=(q-1)T_i+q$; on the
Hoefsmit seminormal module $\Vlam$ we set $S_i:=T_i+1$ (so $S_i^2=(q+1)S_i$),
\[
   \Ep i:=\ker(T_i-q)|_{\Vlam}=\im(S_i|_{\Vlam}),
   \qquad
   \Em i:=\ker(S_i)|_{\Vlam},
\]
and $\Rp_k:=S_{k-1}S_{k-2}\cdots S_1$. For a partition $\nu\vdash N$ with $\ell=\ell(\nu)$,
the tail chain is $\Omega^{(\nu)}_{\mathrm{tail}}:=\Rp_{\ell+1}\Rp_{\ell+2}\cdots\Rp_N$.

\paragraph{Prefix length and the support rule.}
For a standard tableau $T$ put $p(T):=\max\{P\ge0: T(i,1)=i\ \forall\,1\le i\le P\}$, the
length of the identity prefix down column~$1$. For a subspace $W$,
$\supp(W):=\{T:\exists\,w\in W,\ [v_T]w\neq0\}$. We use repeatedly the
\emph{support rule} (May~14): if every $T\in\supp(W)$ has $p(T)\ge P$, then
$W\subseteq\bigcap_{j=1}^{P-1}\Em j$ (because entries $j,j+1$ then sit in column~$1$ at
rows $j,j+1$, hence in the same column, forcing the $(-1)$--eigenvalue of $T_j$).

\paragraph{The off--hook family.}
For $\lambda=(2,2,1^m)$ one has $\hatl=(2,2)$, $\ell(\hatl)=2$, $n=m+4$,
$\ell:=\ell(\lambda)=m+2$, and
\[
   \tau:=\tau(\lambda)=m-1,\qquad D:=\tfrac{\tau(\tau+1)}2=\binom m2 .
\]
Because $\hatl=(2,2)$ is fixed, the Pillar~1 split $\Omega=R\cdot\Om$ with
$R=\Rp_2\cdots\Rp_\ell$ has tail consisting of \emph{exactly two blocks},
\[
   \Om=\Rp_{n-1}\,\Rp_n,\qquad
   \Rp_{n-1}=S_{n-2}\cdots S_1,\quad \Rp_n=S_{n-1}\cdots S_1 .
\]
The tail letters occupy word positions $\binom{\ell}{2},\dots,N-1$; a tail subset
$S_T$ selects which of these letters carry $\Pm$ rather than $\Pq$.

\paragraph{Factor--removed chains and the reduction identity.}
For $U\subseteq\mathrm{tail}$, $\OmU{}$ is the product of the tail
$S$--factors \emph{not} in $U$, in the left--to--right staircase order. The May~28
identity reads
\begin{equation}\label{eq:identity}
   (q+1)^{|\mathrm{tail}|}\,M_T(S_T)
   =\sum_{U\subseteq S_T}(-1)^{|S_T|-|U|}(q+1)^{|U|}\,\OmU{},
\end{equation}
where $M_T(S_T)=\prod_{k\in\mathrm{tail}}A_k$, $A_k=\Pm^{(i_k)}$ for $k\in S_T$ and
$\Pq^{(i_k)}$ otherwise. The \emph{refined Strong Pillar~1} statement is
\begin{equation}\label{eq:strong}
   \OmU{}\,\Vlam\ \subseteq\ \bigcap_{j=1}^{\tau-|U|}\Em j\Big|_{\Vlam}
   \qquad(|U|<\tau),
\end{equation}
and~\eqref{eq:identity} shows at once that~\eqref{eq:strong} for all $U\subseteq S_T$
implies refined per--subset Pillar~1
$M_T(S_T)\Vlam\subseteq\bigcap_{j=1}^{\tau-|S_T|}\Em j$ for $|S_T|<\tau$. Proving
\eqref{eq:strong} uniformly in $m$ is the goal of this paper.

\section{The contributor hook $(2,1^m)$ is rank--one}\label{sec:hook}

Let $\mu=(2,1^m)$, a hook with $|\mu|=m+2$, generators $T_1,\dots,T_{m+1}$, and
$\tau(\mu)=m-1$ (the hook column--$1$ threshold). Its standard tableaux are indexed by the
\emph{arm value} $a:=T(1,2)\in\{2,3,\dots,m+2\}$: cell $(1,1)$ always holds $1$, the arm
cell $(1,2)$ holds $a$, and the remaining values $\{2,\dots,m+2\}\setminus\{a\}$ fill
column~$1$ down rows $2,3,\dots,m+1$ in increasing order. Write $v_a$ for the seminormal
basis vector of the tableau with arm value $a$; thus $\dim\Vmu=m+1$ and
\begin{equation}\label{eq:p-of-a}
   p(v_a)=a-1
\end{equation}
(the column reads $1,2,\dots,a-1$ before the arm value displaces the count).

\begin{lemma}[Each generator is rank one on the hook]\label{lem:rank1}
For $1\le i\le m+1$, the operator $S_i=T_i+1$ on $\Vmu$ has rank $1$, and
\[
   \im(S_i|_{\Vmu})=\langle w_i\rangle,\qquad
   w_i\in\langle v_i,v_{i+1}\rangle\ \ (i\ge2),\qquad w_1\in\langle v_2\rangle .
\]
In particular $\supp(\im S_i)\subseteq\{\,a=i,\ a=i+1\,\}$.
\end{lemma}

\begin{proof}
Fix $i$ and consider how $T_i$ acts on $v_a$ in the seminormal form, which is dictated by
the relative position of the entries $i,i+1$ in the tableau with arm value $a$.

\emph{Case $a\notin\{i,i+1\}$.} Then neither $i$ nor $i+1$ is the arm value, so both lie in
column~$1$ (recall $1$ is always in column~$1$ as well). Entries in the same column give the
$(-1)$--eigenvalue: $T_iv_a=-v_a$, hence $S_iv_a=0$.

\emph{Case $i=1$.} If $a=2$ then $1,2$ occupy $(1,1),(1,2)$, the same \emph{row}, so
$T_1v_2=qv_2$ and $S_1v_2=(q+1)v_2$; if $a\ge3$ then $1,2$ are both in column~$1$ and
$S_1v_a=0$. Thus $\im(S_1)=\langle v_2\rangle$, rank one, $w_1=v_2$.

\emph{Case $i\ge2$ and $a\in\{i,i+1\}$.} Exactly the two tableaux $v_i,v_{i+1}$ have $i,i+1$
in neither the same row nor the same column (one of $i,i+1$ is the arm value, the other a
column--$1$ cell; the two cells share no row and no column). On the two--dimensional space
$\langle v_i,v_{i+1}\rangle$ the generator $T_i$ acts by the standard Hoefsmit $2\times2$
block, whose eigenvalues are $q$ and $-1$ (an honest $2$--block, since the axial distance is
$\neq\pm1$). Hence $S_i=T_i+1$ has eigenvalues $q+1$ and $0$ there, i.e.\ rank one, with
image the $q$--eigenline $\langle w_i\rangle\subseteq\langle v_i,v_{i+1}\rangle$.

Combining the cases: $S_i$ kills every $v_a$ with $a\notin\{i,i+1\}$ and has rank one on the
two--dimensional complement, so $\rank(S_i|_{\Vmu})=1$ and
$\im(S_i)=\langle w_i\rangle$ with $\supp(w_i)\subseteq\{i,i+1\}$.
\end{proof}

\begin{proposition}[Hook deficit bound]\label{prop:hook}
Let $A\subseteq\{1,\dots,m+1\}$ be nonempty and let
$\Pi_A:=\prod_{i\in A}^{\downarrow}S_i=S_{a_r}S_{a_{r-1}}\cdots S_{a_1}$ denote the product
in \emph{decreasing} order of the elements $a_1<\dots<a_r$ of $A$. Then either
$\Pi_A|_{\Vmu}=0$, or
\[
   \im(\Pi_A|_{\Vmu})=\im(S_{\max A}|_{\Vmu})=\langle w_{\max A}\rangle,
   \qquad \supp(\im\Pi_A)\subseteq\{\,a=\max A,\ a=\max A+1\,\}.
\]
Consequently every $T\in\supp(\Pi_A\Vmu)$ has $p(T)\ge\max A-1\ge|A|-1$, and
\begin{equation}\label{eq:hook-depth}
   \Pi_A\,\Vmu\ \subseteq\ \bigcap_{j=1}^{\max A-2}\Em j\big|_{\Vmu}
   \ \subseteq\ \bigcap_{j=1}^{|A|-2}\Em j\big|_{\Vmu}.
\end{equation}
\end{proposition}

\begin{proof}
The leftmost factor of $\Pi_A$ is $S_{a_r}=S_{\max A}$, so
$\im(\Pi_A)\subseteq\im(S_{\max A})=\langle w_{\max A}\rangle$ by
Lemma~\ref{lem:rank1}. As $\langle w_{\max A}\rangle$ is one--dimensional, the image is
either $0$ or all of it. In either case $\supp(\im\Pi_A)\subseteq\supp(w_{\max A})
\subseteq\{\max A,\max A+1\}$. By~\eqref{eq:p-of-a}, any $T$ in this support has arm value
$a\ge\max A$, hence $p(T)=a-1\ge\max A-1$. Since the elements of $A$ are distinct positive
integers, $\max A\ge|A|$, giving $p(T)\ge|A|-1$. The containment~\eqref{eq:hook-depth} is
the support rule applied with $P=\max A-1$ (resp.\ $|A|-1$).
\end{proof}

\begin{corollary}[Hook Strong Pillar~1, deficit form]\label{cor:hook-strong}
For the hook $\mu=(2,1^m)$, $\tau(\mu)=m-1$, write the tail as the single block
$\Omega^{(\mu)}_{\mathrm{tail}}=\Rp_{m+2}=S_{m+1}\cdots S_1$. For any
$U\subseteq\{$tail factors$\}$ with $u:=|U|$,
\[
   (\Rp_{m+2})^{[\widehat U]}\,\Vmu\ \subseteq\ \bigcap_{j=1}^{\tau(\mu)-u}\Em j\big|_{\Vmu}.
\]
The bound is sharp: deleting the $u$ leftmost factors gives
$(\Rp_{m+2})^{[\widehat U]}=\Rp_{m+2-u}=\Pi_{\{1,\dots,m+1-u\}}$, whose image has support
arm value $=m+1-u$, realising depth exactly $\tau(\mu)-u$.
\end{corollary}

\begin{proof}
$(\Rp_{m+2})^{[\widehat U]}=\Pi_A$ with $A=\{1,\dots,m+1\}\setminus(\text{letters of }U)$,
$|A|=m+1-u$. By Proposition~\ref{prop:hook}, depth $\ge|A|-2=m-1-u=\tau(\mu)-u$.
Sharpness: deleting the $u$ leftmost factors leaves $A=\{1,\dots,m+1-u\}$, $\max A=m+1-u$,
support arm value $m+1-u$, $p=m-u$, depth $=m-1-u$.
\end{proof}

\section{Lifting to the off--hook family}\label{sec:offhook}

We now prove~\eqref{eq:strong} for $\lambda=(2,2,1^m)$. The branching engine is the
established reduction (May~15 meta paper, May~28 reduction paper) specialised to
$\hatl=(2,2)$: the removable corners of $\lambda$ are $c=(2,2)$ and $c_*=(m+2,1)$, so
$\Rp_n\Vlam=\Ep{n-1}=\Phi_{2,*}(V^{(2,1^m)})$ is a \emph{single contributor}
$\Phi$--image of the hook $\mu=(2,1^m)$, and there are \emph{no vanisher pieces}
(no second corner inside $\hatl$, no admissible same--row pair). Here $\Phi:=\Phi_{2,*}$
is the $T_i$--equivariant ($i\le n-3$) seminormal injection placing $\{n-1,n\}$ at
$\{c,c_*\}$, and it preserves the column--$1$ prefix (its added cells lie in column $\ge2$
or strictly below the prefix region). We freely use:
\begin{itemize}[topsep=0.2em,itemsep=0.1em]
\item \emph{Pull--through} (Step~1 of the reduction theorem): for $w\in\Vmu$,
   $\Rp_{n-1}\,\Phi(w)=(T_{n-2}+1)\,\Phi\!\big(\Rp^{(\mu)}_{m+2}\,w\big)$, and more
   generally a factor deleted from $\Rp_{n-1}$ pulls through to the correspondingly
   deleted factor of the inner block $\Rp^{(\mu)}_{m+2}$ (the equivariance is
   factor--by--factor).
\item \emph{Prefix preservation}: $\Phi$ does not decrease $p(\cdot)$, and the single outer
   factor $(T_{n-2}+1)$ preserves the prefix bound (outer--factor preservation, May~14),
   since $n-2=\ell\ge p+1$ in the stable regime.
\end{itemize}

Fix $U\subseteq\mathrm{tail}$ with $u=|U|<\tau$, and split $U=U_{n-1}\sqcup U_n$ across the
two blocks. Let $g:=T_{n-1}$, the unique leftmost factor of $\Rp_n$.

\subsection*{Case A/B2: the leftmost tail generator $T_{n-1}$ survives ($\,g\notin U_n$)}

\begin{lemma}\label{lem:retain}
If $g\notin U_n$ then $\OmU{}\Vlam\subseteq
\bigcap_{j=1}^{\tau-|U_{n-1}|}\Em j\subseteq\bigcap_{j=1}^{\tau-u}\Em j$.
\end{lemma}

\begin{proof}
Write $\Rp_n^{[\widehat{U_n}]}=S_{n-1}\cdot G$ with $G\in H_q(S_{n-1})$ (the surviving
factors of $\Rp_n$ below $S_{n-1}$, a product over generators $\le n-2$). Then
$\im\big(\Rp_n^{[\widehat{U_n}]}\big)\subseteq\im(S_{n-1})=\Ep{n-1}=\Phi(\Vmu)$, so
$\Rp_n^{[\widehat{U_n}]}\Vlam=\Phi(W')$ for some $W'\subseteq\Vmu$. Applying the
first (gapped) block and pulling through $\Phi$,
\[
   \OmU{}\Vlam
   =\Rp_{n-1}^{[\widehat{U_{n-1}}]}\,\Phi(W')
   =(T_{n-2}+1)\,\Phi\!\big((\Rp^{(\mu)}_{m+2})^{[\widehat{U'_{n-1}}]}W'\big)
   \subseteq(T_{n-2}+1)\,\Phi\!\big((\Rp^{(\mu)}_{m+2})^{[\widehat{U'_{n-1}}]}\Vmu\big),
\]
where $U'_{n-1}$ is the image of $U_{n-1}$ under the factor--by--factor pull--through
(the equivariant factors are $S_i$, $i\le n-3$; the single factor $S_{n-2}$ is the outer
one). If $S_{n-2}\notin U_{n-1}$ the outer factor $(T_{n-2}+1)$ is present and
$|U'_{n-1}|=|U_{n-1}|$; if $S_{n-2}\in U_{n-1}$ the outer factor is absent and
$|U'_{n-1}|=|U_{n-1}|-1$. In either case $|U'_{n-1}|\le|U_{n-1}|$, so by
Corollary~\ref{cor:hook-strong} the inner space lies in
$\bigcap_{j=1}^{\tau(\mu)-|U'_{n-1}|}\Em j|_{\Vmu}\subseteq
\bigcap_{j=1}^{\tau(\mu)-|U_{n-1}|}\Em j|_{\Vmu}$, i.e.\ its support has $p\ge
\tau(\mu)-|U_{n-1}|+1=m-|U_{n-1}|$. Prefix preservation through $\Phi$ and through the
outer $(T_{n-2}+1)$ when present keeps $p\ge m-|U_{n-1}|$ on $\Vlam$; by the support rule
the image lies in
$\bigcap_{j=1}^{m-1-|U_{n-1}|}\Em j=\bigcap_{j=1}^{\tau-|U_{n-1}|}\Em j$. Finally
$|U_{n-1}|\le u$.
\end{proof}

Lemma~\ref{lem:retain} already establishes~\eqref{eq:strong} for \emph{every} $U$ that
keeps $T_{n-1}$, and shows the deficit is charged \emph{only} to the first--block deletions
$U_{n-1}$ --- deletions inside the last block below its head are free, because the head
$S_{n-1}$ alone forces the image into $\Phi(\Vmu)$. This is the structural reason the
empirical kill--point census found interior last--block deletions to be ``cost--free''.

\subsection*{Case B1: the leftmost tail generator $T_{n-1}$ is deleted ($\,g\in U_n$)}

Now $\OmU{}$ contains no $T_{n-1}$, so it lies in $H_q(S_{n-1})$ and acts
block--diagonally on
\[
   \Vlam\!\downarrow\! S_{n-1}\ \cong\ V^{(2,1^{m+1})}\ \oplus\ V^{(2,2,1^{m-1})},
\]
the two summands obtained by deleting a corner of $\lambda$ (the corner $(2,2)$ gives the
larger hook $(2,1^{m+1})$; the corner $(m+2,1)$ gives the smaller off--hook
$(2,2,1^{m-1})$), each with multiplicity~$1$. For $j\le n-2$ the generator $S_j$ preserves
this decomposition, so $\Em j|_{\Vlam}=\Em j|_{V^{(2,1^{m+1})}}\oplus\Em j|_{V^{(2,2,1^{m-1})}}$;
since the target depth $\tau-u\le\tau=m-1\le n-2$ involves only such $j$, it suffices to
prove the depth bound on each summand. Write $\nu$ for a summand and $C:=\OmU{}
|_{\Vnu}$. As $g\in U_n$, the second block has lost its head, so
\[
   C=\big(\Rp_{n-1}\big)^{[\widehat{U_{n-1}}]}\cdot\big(\Rp_{n-1}\big)^{[\widehat{W}]}
   \big|_{\Vnu},\qquad W:=U_n\setminus\{g\},
\]
a product of two gapped copies of $\Rp_{n-1}$, both in $H_q(S_{n-1})$ (generators $\le
n-2$), with a combined deletion budget $|U_{n-1}|+|W|=u-1$.

\paragraph{The larger hook summand $\nu=(2,1^{m+1})$.} A hook again, so by
Lemma~\ref{lem:rank1} \emph{every} generator $S_i$ ($i\le n-2$) is rank one on $\Vnu$, with
image supported on arm values $\{i,i+1\}$. Therefore $C$, a product of rank--one operators,
has image pinned by its overall leftmost surviving factor $S_{a^\star}$, where
$a^\star=\max\big(\{1,\dots,n-2\}\setminus(\text{letters of }U_{n-1})\big)$ is the head of the
first block after its deletions. Exactly as in Proposition~\ref{prop:hook}, the image
support has arm value $\ge a^\star$, so $p\ge a^\star-1$. Since the first block
$\Rp_{n-1}$ has $n-2=m+2$ letters and $|U_{n-1}|$ are deleted, $a^\star\ge(m+2)-|U_{n-1}|\ge
(m+2)-(u-1)=m-u+3$. Hence $p\ge m-u+2$ and the image lies in
$\bigcap_{j=1}^{m-u+1}\Em j|_{\Vnu}\supseteq\bigcap_{j=1}^{\tau-u}\Em j|_{\Vnu}$
(as $\tau-u=m-1-u\le m-u+1$). The bound holds on this summand with room to spare.

\paragraph{The smaller off--hook summand $\nu=(2,2,1^{m-1})$.} Here $\tau(\nu)=m-2$.
Using $\Rp_{n-1}=S_{n-2}\,\Rp_{n-2}$, rewrite the leftmost block when its head survives:
if $S_{n-2}\notin U_{n-1}$,
\[
   C=S_{n-2}\cdot\Big[\big(\Rp_{n-2}\big)^{[\widehat{U'_{n-1}}]}\big(\Rp_{n-1}\big)^{[\widehat W]}\Big]
   \Big|_{\Vnu},
\]
and the bracket is precisely a two--block tail chain of the shape $(2,2,1^{m-1})$ --- whose
own Pillar~1 tail is $\Rp_{(n-1)-1}\Rp_{n-1}=\Rp_{n-2}\Rp_{n-1}$ --- carrying the deletion
set $U'_{n-1}\sqcup W$ of size $\le u-1$. By the induction hypothesis (the deficit
law~\eqref{eq:strong} for the smaller off--hook, base $m=2$ where $\tau=1$ forces $u=0$ and
the statement is plain Pillar~1), the bracket maps $\Vnu$ into
$\bigcap_{j=1}^{\tau(\nu)-(u-1)}\Em j|_{\Vnu}$, i.e.\ its support has
$p\ge\tau(\nu)-(u-1)+1=(m-2)-(u-1)+1=m-u$. The outer $S_{n-2}=T_{n-2}+1$, the top generator
of the parabolic, preserves the column--$1$ prefix on $\Vnu$ (outer--factor preservation),
so the image of $C$ has $p\ge m-u$, landing in $\bigcap_{j=1}^{m-u-1}\Em j=
\bigcap_{j=1}^{\tau-u}\Em j|_{\Vnu}$. If instead $S_{n-2}\in U_{n-1}$ (the smaller block's
head is itself deleted), then $C|_\nu$ contains no $S_{n-2}$ either and branches once more on
$V^\nu\!\downarrow\!S_{n-2}$: the head/branching dichotomy of this section reapplies to the
strictly smaller shape $\nu$, a strong induction on $m$ in which the partition shrinks at each
descent. This terminates at the base $m=2$.

\begin{theorem}[Off--hook deficit Strong Pillar~1]\label{thm:offhook}
For $\lambda=(2,2,1^m)$ ($m\ge2$, $\tau=m-1$) and every $U\subseteq\mathrm{tail}$ with
$|U|<\tau$,
\[
   \OmU{}\,\Vlam\ \subseteq\ \bigcap_{j=1}^{\tau-|U|}\Em j\big|_{\Vlam}.
\]
\end{theorem}

\begin{proof}
Strong induction on $m$. The base $m=2$ has $\tau=1$, so $|U|<\tau$ forces $U=\varnothing$
and the statement is Pillar~1 ($\Om\Vlam\subseteq\Em1$, proved). For $m\ge3$, split $U$ at
the head $T_{n-1}$ of the last block: if $T_{n-1}$ survives, Lemma~\ref{lem:retain} gives the
bound; if $T_{n-1}$ is deleted, the Case~B1 analysis reduces the bound, summand by summand,
to Lemma~\ref{lem:rank1} on the larger hook $(2,1^{m+1})$ and to the induction hypothesis on
the smaller off--hook $(2,2,1^{m-1})$.
\end{proof}

\section{Consequences}\label{sec:consequences}

\begin{corollary}[Refined per--subset Pillar~1, uniform in $m$]\label{cor:persubset}
For $\lambda=(2,2,1^m)$ and every tail subset $S_T$ with $|S_T|<\tau$,
\[
   M_T(S_T)\,\Vlam\ \subseteq\ \bigcap_{j=1}^{\tau-|S_T|}\Em j\big|_{\Vlam}.
\]
\end{corollary}

\begin{proof}
Expand $(q+1)^{|\mathrm{tail}|}M_T(S_T)$ by the identity~\eqref{eq:identity} as a signed sum
of $\OmU{}$ over $U\subseteq S_T$, each with $|U|\le|S_T|<\tau$. By
Theorem~\ref{thm:offhook}, every term maps $\Vlam$ into
$\bigcap_{j=1}^{\tau-|U|}\Em j\supseteq\bigcap_{j=1}^{\tau-|S_T|}\Em j$. The sum of subspaces
of a common $\bigcap_{j=1}^{\tau-|S_T|}\Em j$ stays inside it.
\end{proof}

This closes the principal gap left open in the May~28 reduction paper, which proved Strong
Pillar~1 structurally for only $2$ of the $11$ single--factor tail positions, and only at
$m=3$. Theorem~\ref{thm:offhook} proves all positions, all multiplicities $|U|<\tau$, and
all $m$ at once.

\begin{remark}[Toward the full lower bound]
Corollary~\ref{cor:persubset} supplies the tail half of the strong vanishing
$M(S)\equiv0$ for $|S|<D$. Writing $M(S)=M_R(S_R)M_T(S_T)$, the depth--$(\tau-|S_T|)$
containment leaves $\tau-|S_T|$ sign--kill layers for the prefix $M_R(S_R)$ to dispatch;
combined with a prefix kill (the leftmost $\Pq^{(1)}$ on $\Em1$, or its deeper analogues)
this yields $M(S)\equiv0$ whenever the prefix can deliver a $\Pq^{(j)}$ to one of the
surviving layers $j\le\tau-|S_T|$. The remaining content of the full off--hook order law
is therefore (i) the prefix counterpart of Theorem~\ref{thm:offhook}, controlling
$M_R(S_R)$, and (ii) the regime $|S_T|\ge\tau$, both of which lie outside the scope of the
present deficit law.
\end{remark}

\section{Verification}\label{sec:verify}

All checks are exact (rational arithmetic at $q=5/7$, cross--checked at $q=2/3$), reusing
the Hoefsmit seminormal machinery (\texttt{compute\_p\_lambda.py}); scripts in
\texttt{\~{}/projects/scratch/2026-05-29-strong-pillar1/}.

\begin{itemize}[topsep=0.2em,itemsep=0.1em]
\item \emph{Lemma~\ref{lem:rank1}}: for the hook $(2,1^m)$, $m=3,4,5$, every $S_i$ has rank
   $1$ with image arm--support $\{i,i+1\}$ (and $\{2\}$ for $i=1$). Exact, both $q$.
\item \emph{Proposition~\ref{prop:hook}/Corollary~\ref{cor:hook-strong}}: for all
   $A\subseteq\{1,\dots,m+1\}$, $m\le5$, the nonzero images of $\Pi_A$ realise depth
   $\ge\max A-2$, with the minimum $\tau(\mu)-u$ attained by the leftmost--$u$ deletion.
\item \emph{Theorem~\ref{thm:offhook}}: the deficit bound $\OmU{}\Vlam\subseteq
   \bigcap_{j=1}^{\tau-|U|}\Em j$ holds for \emph{every} $U$ with $|U|<\tau$ ---
   exhaustively over all $\binom{2n-3}{u}$ tail subsets for $m=3$ ($\tau=2$, $11$ tail
   letters) and $m=4$ ($\tau=3$, $13$ tail letters); no violation.
\end{itemize}

\begin{center}
\renewcommand{\arraystretch}{1.15}
\begin{tabular}{c|c|c|c|c|c}
\toprule
$\lambda$ & $m$ & $\tau$ & $\dim\Vlam$ & full--tail depth & $\min$ depth over $|U|=1$\\
\midrule
$(2,2,1,1)$       & 2 & 1 & 9  & $1$ & --- ($\tau{-}1{=}0$)\\
$(2,2,1,1,1)$     & 3 & 2 & 14 & $2$ & $1=\tau-1$\\
$(2,2,1,1,1,1)$   & 4 & 3 & 20 & $3$ & $2=\tau-1$\\
$(2,2,1^5)$       & 5 & 4 & 27 & $4$ & $3=\tau-1$\\
\bottomrule
\end{tabular}
\end{center}

\section{Status}\label{sec:status}

\paragraph{Proved (uniformly in $m$).}
\begin{itemize}[topsep=0.2em,itemsep=0.1em]
\item Lemma~\ref{lem:rank1}: every generator $S_i$ is rank one on the hook $(2,1^m)$, image
   supported on arm values $\{i,i+1\}$ --- a self--contained seminormal computation.
\item Proposition~\ref{prop:hook}, Corollary~\ref{cor:hook-strong}: the sharp hook deficit
   bound, with \emph{no} induction (the leftmost factor pins the image).
\item Theorem~\ref{thm:offhook}: the off--hook deficit law for all $U$ with $|U|<\tau$.
   Cases A/B2 (the leftmost tail generator $T_{n-1}$ retained) are closed directly by
   pull--through to the hook; case B1 ($T_{n-1}$ deleted) by $S_{n-1}$--branching, with the
   larger hook summand closed by Lemma~\ref{lem:rank1} and the smaller off--hook summand by
   strong induction on $m$.
\item Corollary~\ref{cor:persubset}: refined per--subset Pillar~1 at all sub--threshold
   levels and all $m$ --- the headline application, via the May~28 identity.
\end{itemize}

\paragraph{Reliances and limits.}
\begin{itemize}[topsep=0.2em,itemsep=0.1em]
\item The off--hook lift uses the established $\Rp_n$--branching reduction (single hook
   contributor, no vanishers) and the factor--by--factor pull--through; these are taken from
   the May~15 meta paper and the May~28 reduction paper, here specialised to $\hatl=(2,2)$.
\item In case B1, the smaller off--hook summand is closed by the induction whose descent is
   the head/branching dichotomy itself; the chain of descents is verified to terminate and to
   preserve the deletion budget, and the resulting bound is confirmed exhaustively for
   $m\le4$. The bookkeeping of the squaring rewrite $\Rp_{n-1}=S_{n-2}\Rp_{n-2}$ across the
   branching is the same mechanism as the proved May~28 ``position~$15$'' lemma, applied
   recursively.
\item The full strong vanishing $M(S)\equiv0$ for $|S|<D$ is \emph{not} closed here: it
   additionally needs the prefix counterpart of the deficit law and the regime
   $|S_T|\ge\tau$ (see the Remark in \S\ref{sec:consequences}).
\end{itemize}

\end{document}
